% 10. 窓: すべてを組み合わせる
\section{窓: すべてを組み合わせる}\label{sec:window}

ここまでの結果を組み合わせて、研究の問い(\S\ref{sec:intro})に答える形を作る。

\subsection{二つの条件}

$\mu$ と $f$ を決めると、どの BH がどれだけ回転を失うかが決まる。窓は次の二つを同時に満たす $(\mu,f)$ である。

\paragraph{条件 1: 2 世代目は回転を保つ(上限)}
GW241011 の主星は、生まれたときのスピン $\chi_0$ から、年齢 $T_2$ のあいだに観測の下限 0.69 を下回ってはならない。回転を失う速さは $f^2$ に比例するので(式 \ref{eq:chidot_H})、これは $f$ の上限になる(式 \ref{eq:fmax})。
\begin{equation}
t_\mathrm{2G}(f)\ge T_2\quad\Longleftrightarrow\quad f\le f_\mathrm{max}.
\end{equation}

\paragraph{条件 2: 1 世代目は回転を失う(下限)}
質量とスピンの相関を超放射で説明するには、1 世代目の BH が、年齢 $T_1$ のうちに臨界スピンの近く(臨界 $+\delta$)まで回転を失わなければならない。これは $f$ の下限になる。
\begin{equation}
t_\mathrm{1G}(f)\le T_1\quad\Longleftrightarrow\quad f\ge f_\mathrm{min}(M_1).
\end{equation}
$\delta=0.18$ は、証拠を主張する解析が積み上がりに持たせた幅である。

\paragraph{両立の条件}
$f_\mathrm{min}\le f\le f_\mathrm{max}$ が空でないこと。$t\propto f^{-2}$ なので、これは年齢の比の条件と同じである。
\begin{equation}
\frac{T_1}{T_2}\ge R=\frac{t_\mathrm{1G}}{t_\mathrm{2G}}\qquad(\text{$f$ によらない、\S\ref{sec:equilibrium}}).
\end{equation}
$t\propto1/K\propto\alpha^{-5}$ なので、軽い 1 世代目ほど回転を失うのが遅く、$R$ が大きい。窓は「重い 1 世代目だけが回転を失う」形になりやすい。

\subsection{仮定の組と結果}

\begin{table}[H]
\centering\small
\caption{窓の判定に使う仮定と、$\mu=0.71\times10^{-12}\eV$ での結果(主要項の $\gamma_\mathrm{BH}$。研究ノート E9、E10)。$M_\mathrm{min}$ は 6 Gyr のうちに臨界 $+\delta$ まで回転を失える 1 世代目の最小質量($\delta=0.18/0.05$)。}
\label{tab:window}
\begin{tabularx}{\linewidth}{@{}l>{\raggedright\arraybackslash}X>{\raggedright\arraybackslash}X@{}}
\toprule
 & 楽観的な仮定 & 現実的な仮定 \\
\midrule
GW241011 の主星 & 0.78 で生まれ、年齢 $3\times10^7$ 年 & 0.72 で生まれ、年齢 $5\times10^7$ 年 \\
$f$ の上限(GW241011) & $2.4\times10^{12}\GeV$ & $1.1\times10^{12}\GeV$ \\
$M_\mathrm{min}$(GW だけ) & $8.8/10.9\Msun$ & $11.7/14.5\Msun$ \\
GRS 1915+105(付録) & 年齢 $4.5\times10^7$ 年、0.95 を保つ & 年齢 3 Gyr、0.87 を保つ \\
$M_\mathrm{min}$(GRS も入れる) & --- & $15.5/19.3\Msun$ \\
\bottomrule
\end{tabularx}
\end{table}

\begin{figure}[H]
\centering
% 窓(勉強ノート用)。データ: ../paper/figs/data/fig_window_{opt,real}.dat(研究ノート E9、paper/figs/make_fig_window.py)
\begin{tikzpicture}
\begin{groupplot}[group style={group size=2 by 1, horizontal sep=14mm},
  width=0.48\linewidth, height=0.44\linewidth, note, xmode=log, ymode=log,
  xmin=0.4e-12, xmax=1.2e-12, ymin=3e10, ymax=1e13,
  xtick={0.4e-12,0.6e-12,0.8e-12,1e-12}, xticklabels={0.4,0.6,0.8,1.0},
  xlabel={$\mu\ [10^{-12}\,\mathrm{eV}]$}]
\nextgroupplot[ylabel={$f\ [\mathrm{GeV}]$}, title={\small (a) 楽観的な仮定}]
  \fill[cB!10] (axis cs:0.60e-12,3e10) rectangle (axis cs:0.89e-12,1e13);
  \addplot[cA, line width=1.1pt] table[x=mu, y=GW241011] {../paper/figs/data/fig_window_opt.dat};
  \addplot[ink2, line width=0.8pt, densely dotted] table[x=mu, y=GRS] {../paper/figs/data/fig_window_opt.dat};
  \addplot[cC!60!black, dashed, line width=0.9pt] table[x=mu, y=M15_d18] {../paper/figs/data/fig_window_opt.dat};
  \addplot[cC, dashed, line width=0.9pt] table[x=mu, y=M20_d18] {../paper/figs/data/fig_window_opt.dat};
  \addplot[cC!45, dashed, line width=0.9pt] table[x=mu, y=M25_d18] {../paper/figs/data/fig_window_opt.dat};
  \node[font=\scriptsize, color=cA, anchor=south] at (axis cs:1.0e-12,1.4e12) {GW241011};
  \node[font=\scriptsize, color=cC!60!black, anchor=north east] at (axis cs:1.15e-12,7.5e10) {$15\Msun$};
  \node[font=\scriptsize, color=cC!80!black, anchor=north east] at (axis cs:0.90e-12,4.5e10) {$20\Msun$};
  \node[font=\scriptsize, color=cC!60, anchor=north east] at (axis cs:0.72e-12,4.5e10) {$25\Msun$};
\nextgroupplot[title={\small (b) 現実的な仮定}]
  \fill[cB!10] (axis cs:0.60e-12,3e10) rectangle (axis cs:0.89e-12,1e13);
  \addplot[cA, line width=1.1pt] table[x=mu, y=GW241011] {../paper/figs/data/fig_window_real.dat};
  \addplot[ink2, line width=0.8pt, densely dotted] table[x=mu, y=GRS] {../paper/figs/data/fig_window_real.dat};
  \addplot[cC!60!black, dashed, line width=0.9pt] table[x=mu, y=M15_d18] {../paper/figs/data/fig_window_real.dat};
  \addplot[cC, dashed, line width=0.9pt] table[x=mu, y=M20_d18] {../paper/figs/data/fig_window_real.dat};
  \addplot[cC!45, dashed, line width=0.9pt] table[x=mu, y=M25_d18] {../paper/figs/data/fig_window_real.dat};
  \node[font=\scriptsize, color=cA, anchor=south west] at (axis cs:0.41e-12,4.3e12) {GW241011};
  \node[font=\scriptsize, color=ink2, anchor=north east] at (axis cs:1.18e-12,1.35e11) {GRS(付録)};
\end{groupplot}
\end{tikzpicture}

\caption{$(\mu,f)$ 面の窓。青の実線より下で GW241011 の主星が回転を保つ(条件 1)。緑の破線より上で、その質量の 1 世代目が 6 Gyr のうちに回転を失う(条件 2、$\delta=0.18$)。灰色の点線は付録で扱う GRS 1915+105 の上限。橙の帯は証拠を主張する解析の質量(90\%)。データは論文の図(研究ノート E9)と同じ。}
\label{fig:window_study}
\end{figure}

\subsection{読み方}

\begin{itemize}
\item 楽観的な仮定では、$f\sim10^{12}\GeV$ で 10 $\Msun$ 程度より重い 1 世代目が回転を失い、GW241011 は回転を保てる。
\item 現実的な仮定では、上限が 2 倍ほど下がり、回転を失えるのは 12--15 $\Msun$ 以上の 1 世代目になる。X 線連星 GRS 1915+105 も入れると 15--19 $\Msun$ 以上。
\item 2 世代目の生まれたときのスピンは $\approx0.68$ に集中する(\S\ref{sec:secondgen})。0.78 で生まれる確率は 2\% 以下で、楽観的な仮定はその意味でかなり楽観的である。
\end{itemize}
したがって、最後に決めるのは集団解析である。問いは「GWTC-5 の質量とスピンの相関が、重い 1 世代目だけが回転を失う形を許すか」になる。

\subsection{年齢と判定の幅への依存(研究ノート E12)}

\paragraph{$f$ と年齢は縮退している}
準平衡ではスピンは $\tau=\mu(f/\Mpl)^2T$ だけで決まる(\S\ref{sec:equilibrium})。したがって GW241011 の主星からの上限は、実は $f\sqrt{T_2}$ の上限で、年齢を 4 倍にすると $f$ の上限は半分になる。上の「現実的な」$T_2=5\times10^7$ 年は、$(50+50)\Msun$ の連星の遅れの最小値の目安で、GW241011 の主星の親に近い $(10+10)\Msun$ では $3\times10^8$ 年になる。
\begin{center}\small
\begin{tabular}{rccc}
\toprule
$T_2$ [年] & 生まれたスピン & $f$ の上限 [GeV] & 回転を失う 1 世代目($\chi_0=0.7$、$\delta=0.02/0.05$)\\
\midrule
$5\times10^7$ & 0.72 & $1.1\times10^{12}$ & $16.2/14.5\Msun$ 以上\\
$3\times10^8$ & 0.72 & $4.6\times10^{11}$ & $22.4/20.1\Msun$ 以上\\
$3\times10^8$ & 0.78 & $7.6\times10^{11}$ & $18.8/16.8\Msun$ 以上\\
\bottomrule
\end{tabular}
\end{center}

\paragraph{判定の幅 $\delta$ は 0 にできない}
準平衡では臨界からのずれが $x(t)=[x_0^{-1/2}+Kt/2]^{-2}$ で近づくので(\S\ref{sec:equilibrium})、臨界 $+\delta$ に届く時間は $t=(2/K)(\delta^{-1/2}-x_0^{-1/2})$。$\delta$ が小さいと $t\simeq2/(K\sqrt\delta)$ で、$K\propto f^2$ だから、1 世代目の下限は
\begin{equation}
f_\mathrm{low}\propto\delta^{-1/4}
\end{equation}
で増える。$15\Msun$、6 Gyr、$\chi_0=0.7$ では $\delta=0.18$、0.05、0.01、0.001 で $f_\mathrm{low}=0.56$、1.0、1.7、$3.1\times10^{12}\GeV$。$\delta\to0$ では下限が GW241011 の上限を超え、窓は定義から閉じる。$\delta$ は説明のための判定で、集団解析では置かない。$\chi_\mathrm{SR}$ が予言する臨界からのずれ $x\approx4/(KT)^2\propto\alpha^{-10}f^{-4}T^{-2}$ を、そのまま測定と比べる。

\paragraph{1 世代目の自然スピン}
1 世代目の自然スピンが低い(例: Beta(0.69, 2.28))と、$10$、$15$、$20\Msun$ のうち 211 の臨界より上で生まれるのは 43、30、20\% だけである。残りは 211 の信号をそもそも持たず、一部は 322 だけが育ち、(322, 544) の固定点でゆっくり回転を失う(\S\ref{sec:three_level})。

\subsection{集団解析に渡すもの}

集団解析では、BH ごとに(質量、生まれたときのスピン、年齢)をモンテカルロで振り、$(\mu,f)$ を推定する。そのとき要る関数は、年齢 $T$ のあとのスピン
\begin{equation}
\chi_\mathrm{SR}(M,\chi_0,T;\mu,f)=\chi_T\Bigl(\alpha=GM\mu,\ \chi_0,\ \tau=\mu(f/\Mpl)^2T\Bigr)
\end{equation}
である(\S\ref{sec:equilibrium})。準平衡では三つの量だけの関数なので、表から $10^6$ 回を 1 秒以内で引ける(研究ノート E10--E12)。自由場の側($f$ が大きい)では \S\ref{sec:cloud} の飽和の扱いを使い、両者をつなぐ。

\begin{summary}
\begin{itemize}
\item 窓は「2 世代目は回転を保つ($f$ の上限)」と「1 世代目は回転を失う($f$ の下限)」の重なりで、年齢の比 $T_1/T_2\ge R$ と同じである。
\item 回転を失う速さは $\alpha^5$ に比例するので、窓では重い 1 世代目だけが回転を失う。
\item 現実的な仮定では、窓は $f\sim10^{12}\GeV$ で、回転を失うのは 12--15 $\Msun$ 以上(GW だけ)の 1 世代目。ただし $f$ の上限は $f\sqrt{T_2}$ の上限で、1 世代目の下限は判定の幅 $\delta$ に $\delta^{-1/4}$ で依存する。これを GWTC-5 の集団が許すかを、$\delta$ を使わずに集団解析で決める。
\end{itemize}
\end{summary}
